\documentclass[10pt,a4paper]{amsart}
\usepackage[margin=19mm]{geometry}
\usepackage{amsmath,amssymb,mathtools}
\usepackage[scr=rsfs,cal=euler]{mathalfa}
\usepackage{xfrac}
\usepackage[unicode,hidelinks,bookmarks=false]{hyperref}
\newcommand{\ie}{i.e.\ }
\newcommand{\cf}{cf.\ }
\newcommand{\resp}{respectively\ }
\newcommand{\Subsection}{\subsection}
\providecommand{\nobreakdash}{\nobreak\hbox{-}}
\setlength{\emergencystretch}{2em}
\multlinegap0pt
\usepackage{bm}
\usepackage{bbm}
\usepackage{dsfont}
\usepackage{xspace}
\usepackage{cancel}
\usepackage{mathtools}
\usepackage{amsthm}
\makeatletter
\@ifundefined{theorem}{\newtheorem{theorem}{Theorem}}{}
\@ifundefined{proposition}{\newtheorem{proposition}[theorem]{Proposition}}{}
\makeatother
\makeatletter
\expandafter\ifx\csname ack\endcsname\relax
\newenvironment{ack}
  {\par\medskip\noindent\textbf{Acknowledgements.}\enskip\ignorespaces}{\par}
\fi
\expandafter\ifx\csname funding\endcsname\relax
\newenvironment{funding}
  {\par\medskip\noindent\textbf{Funding.}\enskip\ignorespaces}{\par}
\fi
\newcommand{\mcp}{\mathfrak{I}^{q-\mathrm{Potts}}}
\newcommand{\veco}{\boldsymbol{0}}
\newcommand{\vecs}{\boldsymbol{s}}
\makeatother
\title[A formal framework for higher-order spin models via hypergra]{A formal framework for higher-order spin models via hypergraphs, polymatroids, and the Tutte polynomial}
\author{Khallil Berrekkal, Joanna A. Ellis-Monaghan, Merijn Moody, Cl\'elia de Mulatier}
\date{Source: arXiv:2608.14628v1 (CC BY 4.0)}
\begin{document}
\maketitle

\noindent\emph{Source and licence.} A formal framework for higher-order spin models via hypergraphs, polymatroids, and the Tutte polynomial. Version arXiv:2608.14628v1 (CC BY 4.0), distributed under the Creative Commons Attribution 4.0 International licence, as recorded in the arXiv licence metadata for this exact version and shown on the abstract page https://arxiv.org/abs/2608.14628v1. Full rights notice: licenses/arxiv-2608.14628-CC-BY-4.0.txt.\par\smallskip

\noindent\textit{Editorial scope.} Counted unit: the Proposition whose source label is \texttt{None} in the article (source0.tex lines 1568--1570, source label \texttt{None}), delivered together with the article's own complete original proof of it (source0.tex lines 1571--1614, one \texttt{proof} environment of 44 lines). No mathematics is written, repaired or re-derived: statement and proof are the article's own text line for line. Required context retained uncounted: none. Every definition, display and label of those lines is retained unchanged. Omitted from this package and not redistributed: 1--1567, 1615--1869. Nothing inside a retained excerpt is omitted or altered: every retained body is delivered exactly as the article's lines are written, so no omission receipt of any kind accompanies this package. Citations: the retained excerpts carry no citation command at all, so no bibliography entry of the article is transcribed and no reference is redistributed. Reference adaptation: none was needed and none was made; every reference inside the delivered unit resolves inside it, so no unresolved reference or citation is produced. Printed slips kept literal: no printed word, symbol, hypothesis or number of the counted statement or of its proof was altered. Layout and class: the article's own class is \texttt{amsart} with packages including fontenc, inputenc, amssymb, enumitem, graphicx, booktabs, array, float, xcolor, babel, geometry, hyperref, url, comment; the wrapper is the lane's pinned \texttt{amsart} editorial wrapper at 10pt on A4, which restates the theorem-like environments the retained text needs and adds the article's own macro definitions that the retained text needs (\texttt{\textbackslash mcp}, \texttt{\textbackslash veco}, \texttt{\textbackslash vecs}), with extra wrapper packages (bm, bbm, dsfont, xspace, cancel, mathtools). The delivered numbering is the unit's own. Assets: no retained span contains a figure environment, a graphics-inclusion command, a PDF-page inclusion, a table or a TikZ picture; no image file is copied into this package, the package declares no asset dependency and no image file is redistributed with it. Licence: version arXiv:2608.14628v1 of this article is distributed under the Creative Commons Attribution 4.0 International licence, as recorded in the arXiv licence metadata for this exact version and shown on the abstract page https://arxiv.org/abs/2608.14628v1. A full rights notice is delivered at licenses/arxiv-2608.14628-CC-BY-4.0.txt. Characters: every retained excerpt is pure ASCII, so the delivered engine, XeLaTeX with its default Latin Modern text font, reports no missing character.
\begin{proposition}[Properties of $\mcp$]
    The Delta Potts interaction family $\mcp = (S, \{\Phi^k\}_{k \in \mathbb{Z}_{\geq 0}})$ is group-coset, globally satisfiable, functionally representable and contraction-closed.
\end{proposition}

\begin{proof}
    The Delta Potts interaction family $\mcp$ satisfies the following properties:
    \begin{enumerate}
        \item \textbf{Group-Coset:} The solution set $K_{\Phi^k} = \{ (s, \dots, s) \mid s \in S \}$ corresponds to the diagonal subgroup of the abelian group $S^k$ (where $S$ is viewed as $\mathbb{Z}_q$). Thus, it is a group-coset interaction family.
        \item \textbf{Globally Satisfiable:} The constant configuration $\vecs = \veco$ satisfies $\Phi^k(\veco) = 1$ for all $k$. Thus, the interaction family is globally satisfiable.
        \item \textbf{Functionally Representable:} For $k \in \mathbb{Z}_{>0}$ the constraint $s_i = s_j$ for all $i,j \in [k]$ allows us to express all variables in $[k-1]$ as equal to $s_k$. We set the independent set $I=\{k\}$ and dependent set $D = [k-1]$, with dependency map $\sigma(s_k) = (s_k, \dots, s_k)$.
        \item \textbf{Contraction-closed:} Fix $k \geq 1$, $k' \geq k$, and the
decomposition $I = \{k\}$, $D = [k-1]$, $\sigma(s_k) = (s_k, \dots, s_k)$
above. The extended map $\hat\sigma \colon S^{[k'] \setminus D} \to S^{[k']}$
acts by
\[
    \hat\sigma(\vecs)|_i =
    \begin{cases}
        s_k & i \in [k-1], \\
        s_i & i \in \{k, k+1, \dots, k'\}.
    \end{cases}
\]
In particular, $\hat\sigma(\vecs)|_{[k]}$ is constantly equal to $s_k$. Let
$A \subseteq [k']$; then $\Phi^{|A|}(\hat\sigma(\vecs)|_A) = 1$ if and
only if all coordinates of $\hat\sigma(\vecs)|_A$ are equal. We
distinguish three cases.

\emph{Case 1: $A \cap [k] = \emptyset$.} Then $A \subseteq
\{k+1,\dots,k'\} \subseteq [k'] \setminus D$ and $\hat\sigma(\vecs)|_A =
\vecs|_A$, so $\Phi^{|A|}(\hat\sigma(\vecs)|_A) = \Phi^{|A|}(\vecs|_A)$
and we may take $A' = A$.

\emph{Case 2: $A \subseteq [k]$.} Then $\hat\sigma(\vecs)|_A$ is constantly
$s_k$, so $\Phi^{|A|}(\hat\sigma(\vecs)|_A) = 1$ for all $\vecs$, i.e.\ the
function is constant.

\emph{Case 3: $A \cap [k] \neq \emptyset$ and $A \not\subseteq [k]$.} The
coordinates of $\hat\sigma(\vecs)|_A$ indexed by $A \cap [k]$ are all
equal to $s_k$, while those indexed by $A \cap \{k+1,\dots,k'\}$ equal
$s_i$. All being equal is equivalent to $s_k = s_j$ for every $j \in
A \cap \{k+1,\dots,k'\}$, which is exactly
$\Phi^{|A'|}(\vecs|_{A'})$ for $A' := \{k\} \cup (A \cap
\{k+1,\dots,k'\}) \subseteq [k'] \setminus D$.

In each case the function is either constant or equals
$\Phi^{|A'|}(\vecs|_{A'})$ for some unique $A' \subseteq [k'] \setminus D$,
verifying contraction-closedness. \qedhere
\end{enumerate}
\end{proof}

\end{document}
